\documentclass[12pt]{article}

% ============================================================
% Packages
% ============================================================
\usepackage[a4paper, margin=2.5cm]{geometry}
\usepackage{amsmath, amssymb, amsfonts}
\usepackage{graphicx}
\usepackage{booktabs}
\usepackage{hyperref}
\usepackage{cleveref}
\usepackage{float}
\usepackage{subcaption}
\usepackage{algorithm}
\usepackage{algpseudocode}
\usepackage{bm}
\usepackage{natbib}
\usepackage{xcolor}

\hypersetup{
    colorlinks=true,
    linkcolor=blue,
    citecolor=blue,
    urlcolor=blue
}

% ============================================================
% Custom commands
% ============================================================
\newcommand{\R}{\mathbb{R}}
\newcommand{\norm}[1]{\left\| #1 \right\|}
\newcommand{\abs}[1]{\left| #1 \right|}
\newcommand{\tensor}[1]{\boldsymbol{\mathcal{#1}}}
\newcommand{\mat}[1]{\mathbf{#1}}
\newcommand{\vect}[1]{\mathbf{#1}}
\newcommand{\dt}{\Delta t}

% ============================================================
\title{Tensor Train Decomposition with Operator Inference for Parametric Wind Turbine Wake Flow Prediction}

\author{
  Mandar Tabib \\
  \textit{SINTEF Digital, Department of Mathematics and Cybernetics} \\
  \textit{Trondheim, Norway} \\
  mandar.tabib@sintef.no
}

\date{}

% ============================================================
\begin{document}
\maketitle

% ============================================================
% ABSTRACT
% ============================================================
\begin{abstract}
High-fidelity computational fluid dynamics (CFD) simulations of wind turbine wakes are essential for optimising yaw-based wake steering strategies, yet their computational cost prohibits real-time deployment.
This paper presents a data-driven surrogate model that combines Tensor Train (TT) decomposition, Operator Inference (OpInf), and radial basis function (RBF) interpolation---collectively termed \emph{TT-OpInf}---for parametric, spatio-temporal prediction of turbine wake velocity fields under varying yaw angles.
The TT decomposition separates the four-dimensional snapshot tensor (parameter $\times$ time $\times$ space $\times$ velocity component) into three compact cores representing parameter, temporal, and spatial modes, achieving a $10.1\times$ data compression.
OpInf learns a linear dynamical system governing the temporal evolution of the reduced coefficients, while thin-plate-spline RBF interpolation enables prediction at unseen yaw angles.
Trained on OpenFOAM simulations at three yaw angles (270\textdegree, 282\textdegree, 285\textdegree) on an unstructured mesh with 180{,}857 nodes and 101 time steps, the TT-OpInf model predicts the three-component velocity field at a held-out yaw angle of 276\textdegree{} with a relative $L^2$ error of 12.9\% and sub-second inference time.
The method requires no neural network training, offering an interpretable and computationally lightweight alternative to deep-learning-based surrogates.
\end{abstract}

\noindent\textbf{Keywords:} reduced-order model, tensor train decomposition, operator inference, wind turbine wake, yaw steering, surrogate model

% ============================================================
% 1. INTRODUCTION
% ============================================================
\section{Introduction}
\label{sec:introduction}

Wind energy has become a cornerstone of the global transition to renewable power, yet wake-induced losses remain a significant challenge for wind farm operators.
Downstream turbines in a farm can experience power deficits of 10--20\% due to the reduced wind speed and increased turbulence within the upstream turbine's wake \citep{Barthelmie2009,PorteAgel2020}.
Intentional yaw misalignment of upstream turbines---commonly referred to as \emph{wake steering}---deflects the wake laterally away from downstream rotors, thereby increasing aggregate farm power output \citep{Fleming2017,Gebraad2016}.
Effective implementation of wake steering requires fast, accurate predictions of the three-dimensional wake velocity field as a function of yaw angle.

High-fidelity CFD solvers such as OpenFOAM \citep{OpenFOAM} can provide detailed wake flow fields for any yaw configuration; however, a single simulation with adequate spatio-temporal resolution can require hours to days on parallel computing clusters.
This computational burden makes direct CFD impractical for real-time control applications, design optimisation, or uncertainty quantification, all of which demand thousands of evaluations across the parameter space.

To bridge this gap, \emph{reduced-order models} (ROMs) and data-driven surrogates have attracted growing attention.
Classical approaches based on Proper Orthogonal Decomposition (POD) and Dynamic Mode Decomposition (DMD) project the high-dimensional state onto a small set of spatial modes and learn their temporal evolution \citep{Brunton2019,Benner2015}.
More recently, deep-learning surrogates---including convolutional autoencoders \citep{Tabib2024}, graph neural networks \citep{Chen2024}, and neural operators \citep{Li2021fourier}---have demonstrated the ability to learn complex spatio-temporal dynamics directly from data.
While powerful, these methods often require large training datasets, extensive hyperparameter tuning, and GPU resources, and they can lack the interpretability of classical ROM techniques.

An alternative direction is offered by \emph{tensor decomposition} methods, which exploit the multi-dimensional structure of parametric snapshot data.
The Tensor Train (TT) format \citep{Oseledets2011} provides a compact, hierarchical decomposition that separates parameter, temporal, and spatial modes with controllable accuracy.
When combined with \emph{Operator Inference} (OpInf) \citep{Peherstorfer2016}, which learns polynomial dynamical systems from snapshot data via least-squares regression, the result is a framework that captures temporal dynamics without the overhead of neural network training.

In this work, we propose the \emph{TT-OpInf} framework for parametric wind turbine wake prediction.
The key contributions are:
\begin{enumerate}
    \item A three-core TT decomposition of the four-dimensional CFD snapshot tensor that separates parameter, temporal, and spatial variability, achieving $10.1\times$ compression with 2.6\% reconstruction error.
    \item Integration of linear OpInf for temporal dynamics and thin-plate-spline RBF interpolation for parameter generalisation, enabling prediction at unseen yaw angles.
    \item Demonstration on OpenFOAM wake data with 180{,}857 spatial nodes, showing 12.9\% relative $L^2$ error at a held-out yaw angle with sub-second inference time.
\end{enumerate}

The remainder of the paper is organised as follows.
\Cref{sec:methodology} describes the TT-OpInf methodology.
\Cref{sec:data} details the CFD dataset and training/test configuration.
\Cref{sec:results} presents training and inference results.
\Cref{sec:conclusion} summarises findings and outlines future work.


% ============================================================
% 2. METHODOLOGY
% ============================================================
\section{Methodology}
\label{sec:methodology}

\subsection{Problem formulation}
\label{sec:problem}

Consider a set of time-dependent CFD simulations of wind turbine wakes, each characterised by a yaw angle parameter $\mu \in \{\mu_1, \dots, \mu_P\}$.
At each parameter value, the solver produces a sequence of velocity snapshots $\vect{u}(\mu, t_k, \vect{x}) \in \R^3$ for $k = 0, 1, \dots, K$ on an unstructured spatial mesh with $M$ nodes.
Collecting all snapshots yields a four-dimensional data tensor:
\begin{equation}
\label{eq:tensor4d}
    \tensor{T} \in \R^{P \times (K+1) \times M \times 3}, \qquad
    \tensor{T}_{p,k,m,c} = u_c(\mu_p, t_k, \vect{x}_m),
\end{equation}
where index $p$ runs over parameters, $k$ over time steps, $m$ over spatial nodes, and $c$ over velocity components $(U_x, U_y, U_z)$.

The goal is to construct a surrogate model $\hat{\vect{u}}(\mu_{\text{new}}, t_k, \vect{x}_m)$ that predicts the full spatio-temporal velocity field at a new yaw angle $\mu_{\text{new}}$ not present in the training set, at a fraction of the CFD computational cost.

\subsection{Data normalisation}
\label{sec:normalisation}

Prior to decomposition, the tensor is normalised element-wise to $[0,1]$:
\begin{equation}
    \tilde{\tensor{T}} = \frac{\tensor{T} - T_{\min}}{T_{\max} - T_{\min}},
\end{equation}
where $T_{\min}$ and $T_{\max}$ are the global minimum and maximum values across the entire tensor.
Predictions are denormalised by the inverse transformation after reconstruction.

\subsection{Tensor Train decomposition}
\label{sec:tt}

The TT decomposition \citep{Oseledets2011} factorises the normalised data tensor into a chain of low-rank cores.
We first flatten the spatial and feature dimensions to form a three-dimensional tensor $\tilde{\tensor{T}}^{(3)} \in \R^{P \times (K+1) \times (M \cdot 3)}$, then decompose it as
\begin{equation}
\label{eq:tt}
    \tilde{\tensor{T}}^{(3)}_{p,k,s} \approx \sum_{i=1}^{r_1} \sum_{j=1}^{r_2} G^{(1)}_{p,i} \; G^{(2)}_{i,k,j} \; G^{(3)}_{j,s},
\end{equation}
where:
\begin{itemize}
    \item $\mat{G}^{(1)} \in \R^{P \times r_1}$ is the \emph{parameter core}, encoding dependence on yaw angle;
    \item $\mat{G}^{(2)} \in \R^{r_1 \times (K+1) \times r_2}$ is the \emph{time core}, encoding temporal dynamics;
    \item $\mat{G}^{(3)} \in \R^{r_2 \times (M \cdot 3)}$ is the \emph{spatial core}, encoding spatial modes.
\end{itemize}

The decomposition is computed by two successive truncated SVDs.
First, the tensor is unfolded along the parameter mode to form $\mat{X}_1 \in \R^{P \times (K+1) \cdot M \cdot 3}$ and factorised as
\begin{equation}
    \mat{X}_1 = \mat{U}_1 \mat{\Sigma}_1 \mat{V}_1^\top \approx \mat{U}_1^{(r_1)} \mat{\Sigma}_1^{(r_1)} {\mat{V}_1^{(r_1)}}^\top,
\end{equation}
where the rank $r_1$ is determined by the truncation tolerance $\varepsilon$ such that
\begin{equation}
    \frac{\sum_{i=r_1+1}^{\min(P,N)} \sigma_i^2}{\sum_{i=1}^{\min(P,N)} \sigma_i^2} \leq \varepsilon^2.
\end{equation}

Setting $\mat{G}^{(1)} = \mat{U}_1^{(r_1)}$, the intermediate matrix $\mat{\Sigma}_1^{(r_1)} {\mat{V}_1^{(r_1)}}^\top$ is reshaped to $(r_1 \cdot (K+1)) \times (M \cdot 3)$ and a second truncated SVD yields $\mat{G}^{(2)}$ and $\mat{G}^{(3)}$ with rank $r_2$.

Using Einstein summation notation, the full reconstruction reads:
\begin{equation}
\label{eq:reconstruct}
    \hat{\tilde{\tensor{T}}}^{(3)}_{p,k,s} = \mat{G}^{(1)}_{p,:} \; \mat{G}^{(2)}_{:,k,:} \; \mat{G}^{(3)}_{:,s}.
\end{equation}

The \emph{compression ratio} is defined as
\begin{equation}
    \rho = \frac{P \cdot (K+1) \cdot M \cdot 3}{P \cdot r_1 + r_1 \cdot (K+1) \cdot r_2 + r_2 \cdot M \cdot 3}.
\end{equation}

\subsection{Operator Inference for temporal dynamics}
\label{sec:opinf}

To enable prediction at future time steps or under modified initial conditions, we learn a dynamical system governing the temporal evolution of the reduced time coefficients.
From the time core $\mat{G}^{(2)} \in \R^{r_1 \times (K+1) \times r_2}$, we extract a matrix of time coefficients:
\begin{equation}
    \vect{g}(t_k) \in \R^{r_1 \cdot r_2}, \qquad
    \vect{g}(t_k) = \text{vec}\!\left(\mat{G}^{(2)}_{:,k,:}\right),
\end{equation}
obtained by transposing and reshaping the time core to $\mat{G}_{\text{time}} \in \R^{(K+1) \times (r_1 \cdot r_2)}$.

Following the OpInf framework \citep{Peherstorfer2016}, we posit a linear dynamical system:
\begin{equation}
\label{eq:opinf}
    \frac{d\vect{g}}{dt} = \vect{c} + \mat{A}\,\vect{g},
\end{equation}
where $\vect{c} \in \R^{r_1 r_2}$ is a constant vector and $\mat{A} \in \R^{(r_1 r_2) \times (r_1 r_2)}$ is the system matrix.

Time derivatives are approximated by central finite differences at interior points:
\begin{equation}
    \dot{\vect{g}}(t_k) \approx \frac{\vect{g}(t_{k+1}) - \vect{g}(t_{k-1})}{2\,\dt}, \qquad k = 1, \dots, K-1.
\end{equation}

Assembling the data matrix $\mat{D} = [\vect{1}, \mat{G}_{\text{time}}^{(1:K-1)}] \in \R^{(K-1) \times (1 + r_1 r_2)}$ and the derivative matrix $\dot{\mat{G}} \in \R^{(K-1) \times r_1 r_2}$, the operators are obtained by Tikhonov-regularised least squares:
\begin{equation}
\label{eq:opinf_solve}
    [\vect{c}^\top; \mat{A}^\top] = \left(\mat{D}^\top \mat{D} + \lambda \mat{I}\right)^{-1} \mat{D}^\top \dot{\mat{G}},
\end{equation}
where $\lambda > 0$ is the regularisation parameter.

During inference, given the initial condition $\vect{g}_0 = \vect{g}(t_0)$, the trajectory $\{\vect{g}(t_k)\}_{k=0}^{K}$ is computed via fourth-order Runge--Kutta (RK4) integration:
\begin{equation}
    \vect{g}_{k+1} = \vect{g}_k + \frac{\dt}{6}\left(\vect{k}_1 + 2\vect{k}_2 + 2\vect{k}_3 + \vect{k}_4\right),
\end{equation}
where $\vect{k}_i$ are the standard RK4 intermediate evaluations of the right-hand side of \cref{eq:opinf}.

\subsection{RBF interpolation for new parameters}
\label{sec:rbf}

To predict at a new yaw angle $\mu_{\text{new}}$ not in the training set, we interpolate the parameter core $\mat{G}^{(1)}$ using thin-plate-spline radial basis functions.
Given training parameters $\{\mu_1, \dots, \mu_P\}$ and corresponding rows of $\mat{G}^{(1)}$, the RBF interpolant yields:
\begin{equation}
    \hat{\vect{g}}^{(1)}(\mu_{\text{new}}) = \sum_{p=1}^{P} w_p \, \phi\!\left(\norm{\mu_{\text{new}} - \mu_p}\right),
\end{equation}
where $\phi(r) = r^2 \ln r$ is the thin-plate spline kernel and $\{w_p\}$ are weights determined by the interpolation conditions.

\subsection{Inference pipeline}
\label{sec:pipeline}

The complete prediction pipeline for a new yaw angle $\mu_{\text{new}}$ proceeds as follows:

\begin{algorithm}[H]
\caption{TT-OpInf Inference}
\label{alg:inference}
\begin{algorithmic}[1]
\Require New yaw angle $\mu_{\text{new}}$, number of time steps $K$, time step $\dt$
\Ensure Predicted velocity field $\hat{\vect{u}} \in \R^{(K+1) \times M \times 3}$
\State $\hat{\vect{g}}^{(1)} \gets \text{RBF}(\mu_{\text{new}}; \{\mu_p, \mat{G}^{(1)}_p\}_{p=1}^{P})$
\Comment{Interpolate parameter core}
\State $\vect{g}_0 \gets \mat{G}_{\text{time}}[0, :]$
\Comment{Initial time coefficient}
\State $\{\vect{g}_k\}_{k=0}^{K} \gets \text{RK4}(\vect{g}_0, \vect{c}, \mat{A}, K, \dt)$
\Comment{Integrate OpInf dynamics}
\For{$k = 0, \dots, K$}
    \State $\hat{\tilde{\vect{u}}}(t_k) \gets \hat{\vect{g}}^{(1)} \cdot \text{reshape}(\vect{g}_k) \cdot \mat{G}^{(3)}$
    \Comment{TT reconstruction}
\EndFor
\State $\hat{\vect{u}} \gets \text{denormalise}(\hat{\tilde{\vect{u}}})$
\Comment{Map back to physical units}
\end{algorithmic}
\end{algorithm}


% ============================================================
% 3. DATA DESCRIPTION
% ============================================================
\section{CFD Dataset}
\label{sec:data}

\subsection{Simulation setup}

The training data are generated using OpenFOAM \citep{OpenFOAM}, a widely used open-source CFD framework, to simulate wind turbine wake flows under varying yaw misalignment angles.
The computational domain is clipped to the wake-affected region downstream of the turbine rotor, spanning $[960, 2710]$\,m in the streamwise ($x$) direction and $[1190, 1810]$\,m in the lateral ($y$) direction, at a fixed hub height of $z = 88$\,m.
The unstructured mesh contains $M = 180{,}857$ nodes with $1{,}076{,}982$ undirected edges.

The solver outputs the three-component velocity vector $\vect{u} = (U_x, U_y, U_z)$ at each mesh node for 101 equally spaced time steps with $\dt = 0.1$\,s, corresponding to a total simulation window of 10.0\,s.
Velocity magnitudes across the dataset range from 2.74\,m/s to 9.88\,m/s with a mean of 2.34\,m/s.

\subsection{Parametric dataset}

Four yaw angles are simulated: $\mu \in \{270\degree, 276\degree, 282\degree, 285\degree\}$.
The complete four-dimensional snapshot tensor has dimensions:
\begin{equation}
    \tensor{T} \in \R^{4 \times 101 \times 180{,}857 \times 3},
\end{equation}
totalling approximately 220 million scalar values.

\subsection{Training and test split}

The dataset is split by \emph{parameter}, holding out the yaw angle $\mu = 276\degree$ for testing.
This tests the model's ability to \emph{interpolate} to a parameter value lying between training samples:
\begin{itemize}
    \item \textbf{Training:} $\mu_{\text{train}} = \{270\degree, 282\degree, 285\degree\}$ --- 3 yaw angles, all 101 time steps;
    \item \textbf{Test:} $\mu_{\text{test}} = 276\degree$ --- 1 yaw angle (held out), all 101 time steps.
\end{itemize}

The training tensor thus has shape $(3, 101, 180{,}857, 3)$ and the test tensor has shape $(1, 101, 180{,}857, 3)$.


% ============================================================
% 4. RESULTS AND DISCUSSION
% ============================================================
\section{Results and Discussion}
\label{sec:results}

\subsection{TT decomposition analysis}
\label{sec:tt_results}

\Cref{tab:tt_config} summarises the TT-OpInf hyperparameters used for training.
The TT decomposition of the training tensor $(3, 101, 180{,}857, 3)$ produces three cores with ranks $r_1 = 3$ and $r_2 = 30$:
\begin{itemize}
    \item Parameter core: $\mat{G}^{(1)} \in \R^{3 \times 3}$,
    \item Time core: $\mat{G}^{(2)} \in \R^{3 \times 101 \times 30}$,
    \item Spatial core: $\mat{G}^{(3)} \in \R^{30 \times 542{,}571}$.
\end{itemize}

The total number of stored parameters is $9 + 9{,}090 + 16{,}277{,}130 \approx 16.3 \times 10^6$, compared to $164.4 \times 10^6$ in the original tensor, yielding a compression ratio of $\rho = 10.1\times$.
The relative $L^2$ reconstruction error on the training data is $2.59 \times 10^{-2}$ (2.6\%), confirming that the TT decomposition captures the dominant spatial, temporal, and parametric variability.

The first rank $r_1 = 3$ equals the number of training parameters, indicating that the parameter mode retains full information.
The second rank $r_2 = 30$ reflects the intrinsic temporal complexity of the wake flow.

\begin{table}[t]
\centering
\caption{TT-OpInf configuration and training summary.}
\label{tab:tt_config}
\begin{tabular}{lcc}
\toprule
\textbf{Component} & \textbf{Parameter} & \textbf{Value} \\
\midrule
TT decomposition & Tolerance $\varepsilon$ & $10^{-4}$ \\
                 & Max rank & 30 \\
                 & Rank $r_1$ (parameter) & 3 \\
                 & Rank $r_2$ (time) & 30 \\
                 & Reconstruction error & 2.59\% \\
                 & Compression ratio & 10.1$\times$ \\
\midrule
Operator Inference & Polynomial degree & 1 (linear) \\
                   & Regularisation $\lambda$ & $10^{-3}$ \\
                   & System dimension & 90 ($= r_1 \times r_2$) \\
                   & Time integrator & RK4 \\
                   & Fit error & 0.437 \\
\midrule
RBF interpolation & Kernel & Thin-plate spline \\
\midrule
Normalisation & Range & $[0, 1]$ (min-max) \\
\midrule
Training time & Total & $\sim$5 seconds (CPU) \\
Inference time & Per yaw angle & $<$1 second \\
\bottomrule
\end{tabular}
\end{table}

\subsection{OpInf dynamics training}

The OpInf model learns a linear system (\cref{eq:opinf}) of dimension $r_1 \times r_2 = 90$.
The least-squares fit error is $4.37 \times 10^{-1}$, which reflects the challenge of capturing the full temporal variability in a purely linear model.
Nonetheless, the learned dynamics combined with TT reconstruction yield accurate predictions, as shown below.
A higher-order (quadratic) OpInf model was also tested but did not improve accuracy for this dataset, likely due to the limited number of training parameters.

The total offline training time---including data normalisation, TT decomposition, and OpInf fitting---is approximately 5 seconds on a standard CPU, making the method orders of magnitude faster than neural-network-based alternatives that require GPU training.

\subsection{Prediction at held-out yaw angle}
\label{sec:test_results}

The trained TT-OpInf model is used to predict the velocity field at the held-out yaw angle $\mu = 276\degree$ for 100 time steps.
The three-stage inference pipeline---RBF interpolation of $\mat{G}^{(1)}$, RK4 integration of OpInf dynamics, and TT reconstruction---completes in under one second.

\subsubsection{Quantitative error analysis}

\Cref{tab:error_summary} reports the prediction error metrics.
The overall relative $L^2$ error is $12.92\%$.
Importantly, the per-timestep error remains stable across the entire prediction horizon, ranging from $12.2\%$ to $13.6\%$ with a mean of $12.9\%$, indicating no error accumulation from the OpInf time integration.

\begin{table}[t]
\centering
\caption{Prediction error at the held-out yaw angle $\mu = 276\degree$.}
\label{tab:error_summary}
\begin{tabular}{lc}
\toprule
\textbf{Metric} & \textbf{Value} \\
\midrule
Overall relative $L^2$ error & 12.92\% \\
Mean per-timestep error & 12.90\% \\
Maximum per-timestep error & 13.57\% \\
Minimum per-timestep error & 12.20\% \\
\midrule
Interpolation-only error (no OpInf) & 13.05\% \\
OpInf improvement over interpolation & 0.13\% \\
\bottomrule
\end{tabular}
\end{table}

To isolate the contribution of OpInf, we also evaluate a baseline that uses only RBF-interpolated TT reconstruction without temporal dynamics learning.
This interpolation-only baseline achieves $13.05\%$ error, compared to $12.92\%$ with OpInf, confirming that the learned dynamics provide a modest but consistent improvement.
The relatively small margin reflects the dominance of the spatial TT truncation error ($2.6\%$) over the temporal dynamics error.

\subsubsection{Temporal error evolution}

\Cref{fig:error_time} shows the relative $L^2$ error as a function of time step.
The error oscillates mildly between 12.2\% and 13.6\% without a systematic upward trend, demonstrating the temporal stability of the RK4 integration of the learned OpInf system.
The cyclic pattern correlates with the transient wake dynamics in the CFD data.

\begin{figure}[t]
    \centering
    \includegraphics[width=0.9\textwidth]{../tt_opinf_error_over_time_276.png}
    \caption{Relative $L^2$ prediction error over 100 time steps for the held-out yaw angle $\mu = 276\degree$. The error remains stable between 12.2\% and 13.6\% with no accumulation.}
    \label{fig:error_time}
\end{figure}

\subsubsection{Spatial comparison}

\Cref{fig:comparison} presents velocity magnitude contour plots comparing ground truth CFD, TT-OpInf predictions, and absolute error maps at five representative time steps ($t = 0, 25, 50, 75, 100$).
The surrogate model captures the overall wake structure---including the velocity deficit region, lateral wake deflection, and wake recovery---with high fidelity.
The maximum absolute error at any single time step is approximately 3.4\,m/s, localised at high-gradient regions near the wake boundaries and turbulent eddies, where the low-rank TT approximation smooths fine-scale features.

\begin{figure}[t]
    \centering
    \includegraphics[width=\textwidth]{../tt_opinf_comparison_276.png}
    \caption{Velocity magnitude comparison at yaw = 276\textdegree{} for $t = 0, 25, 50, 75, 100$. Left: ground truth (CFD). Centre: TT-OpInf prediction. Right: absolute error. Maximum absolute errors are annotated on each panel.}
    \label{fig:comparison}
\end{figure}

\subsection{Discussion}

The TT-OpInf framework achieves several desirable properties for wind turbine wake surrogates:

\paragraph{Computational efficiency.}
The offline training cost ($\sim$5 seconds) and online inference cost ($<$1 second per yaw angle for 101 time steps on 180{,}857 nodes) are negligible compared to the hours required by CFD.
This speedup factor of $\mathcal{O}(10^3\text{--}10^4)$ enables applications in real-time control and large-scale optimisation.

\paragraph{Interpretability.}
The three TT cores provide physical insight: $\mat{G}^{(1)}$ captures the parametric (yaw-angle) dependence, $\mat{G}^{(2)}$ encodes the dominant temporal modes, and $\mat{G}^{(3)}$ represents spatial patterns.
This separation is not available in black-box neural network surrogates.

\paragraph{No neural network training.}
Unlike GCA-ROM \citep{Chen2024} or DeepONet \citep{Lu2021deeponet} based approaches, TT-OpInf requires no GPU, no iterative gradient-based optimisation, and no hyperparameter searches for learning rate, batch size, or network architecture.

\paragraph{Error sources.}
The prediction error ($\sim$13\%) is dominated by the TT truncation error (2.6\%), which arises from the limited rank approximation.
The spatial core $\mat{G}^{(3)}$ smooths fine-scale turbulent structures, particularly at wake boundaries.
Increasing $r_2$ beyond 30 would reduce this error at the cost of larger storage and potential overfitting with only three training parameters.
The OpInf dynamics contribute minimally to the total error (0.13\% improvement over static interpolation), suggesting that the temporal variability is well captured by the TT time core.

\paragraph{Limitations.}
The current framework is limited to \emph{interpolation} within the training parameter range; extrapolation to yaw angles outside $[270\degree, 285\degree]$ was not tested.
The linear OpInf model may not capture strongly nonlinear temporal dynamics in more turbulent regimes.
The accuracy could be improved by incorporating more training yaw angles or by adopting a multi-fidelity approach where a neural network learns residual corrections to the TT-OpInf predictions.


% ============================================================
% 5. CONCLUSION
% ============================================================
\section{Conclusion}
\label{sec:conclusion}

This paper presented TT-OpInf, a data-driven surrogate framework combining Tensor Train decomposition, Operator Inference, and RBF interpolation for parametric prediction of wind turbine wake velocity fields under varying yaw angles.
The method achieves a 10.1$\times$ compression of the CFD snapshot tensor, trains in approximately 5 seconds on a standard CPU, and predicts the full three-component velocity field on 180{,}857 mesh nodes at sub-second speed.
For a held-out yaw angle of 276\textdegree, the model attains a relative $L^2$ error of 12.9\%, with stable error across 100 predicted time steps.

The TT-OpInf approach offers an interpretable, lightweight alternative to deep-learning surrogates, particularly attractive when training data are limited and rapid deployment is required.
Future work will explore: (i) multi-fidelity extensions where a neural network learns residual corrections to improve accuracy, (ii) inclusion of additional parameters (wind speed, turbulence intensity) beyond yaw angle, and (iii) extension to full three-dimensional wake fields across multiple hub heights.


% ============================================================
% ACKNOWLEDGEMENTS
% ============================================================
\section*{Acknowledgements}
This work was supported by SINTEF Digital. The authors acknowledge computational resources provided by [to be added].


% ============================================================
% REFERENCES
% ============================================================
\bibliographystyle{plainnat}

\begin{thebibliography}{99}

\bibitem[Barthelmie et~al.(2009)]{Barthelmie2009}
R.J.~Barthelmie, K.~Hansen, S.T.~Frandsen, O.~Rathmann, J.G.~Schepers, W.~Schlez, J.~Phillips, K.~Rados, A.~Zervos, E.S.~Politis, and P.K.~Chaviaropoulos.
\newblock Modelling and measuring flow and wind turbine wakes in large wind farms offshore.
\newblock \emph{Wind Energy}, 12(5):431--444, 2009.

\bibitem[Benner et~al.(2015)]{Benner2015}
P.~Benner, S.~Gugercin, and K.~Willcox.
\newblock A survey of projection-based model reduction methods for parametric dynamical systems.
\newblock \emph{SIAM Review}, 57(4):483--531, 2015.

\bibitem[Brunton and Kutz(2019)]{Brunton2019}
S.L.~Brunton and J.N.~Kutz.
\newblock \emph{Data-Driven Science and Engineering: Machine Learning, Dynamical Systems, and Control}.
\newblock Cambridge University Press, 2019.

\bibitem[Chen et~al.(2024)]{Chen2024}
J.~Chen, M.~Tabib, and others.
\newblock Time extrapolation with graph convolutional autoencoder and tensor train decomposition.
\newblock \emph{arXiv preprint arXiv:2511.23037}, 2024.

\bibitem[Fleming et~al.(2017)]{Fleming2017}
P.~Fleming, P.M.O.~Gebraad, S.~Lee, J.-W.~van~Wingerden, K.~Johnson, M.~Churchfield, J.~Michalakes, P.~Spalart, and P.~Moriarty.
\newblock Simulation comparison of wake mitigation control strategies for a two-turbine case.
\newblock \emph{Wind Energy}, 18(12):2135--2143, 2017.

\bibitem[Gebraad et~al.(2016)]{Gebraad2016}
P.M.O.~Gebraad, F.W.~Teeuwisse, J.W.~van~Wingerden, P.A.~Fleming, S.D.~Ruben, J.R.~Marden, and L.Y.~Pao.
\newblock Wind plant power optimization through yaw control using a parametric model for wake effects---a {CFD} simulation study.
\newblock \emph{Wind Energy}, 19(1):95--114, 2016.

\bibitem[Li et~al.(2021)]{Li2021fourier}
Z.~Li, N.~Kovachki, K.~Azizzadenesheli, B.~Liu, K.~Bhatt, A.~Stuart, and A.~Anandkumar.
\newblock Fourier neural operator for parametric partial differential equations.
\newblock In \emph{International Conference on Learning Representations (ICLR)}, 2021.

\bibitem[Lu et~al.(2021)]{Lu2021deeponet}
L.~Lu, P.~Jin, G.~Pang, Z.~Zhang, and G.E.~Karniadakis.
\newblock Learning nonlinear operators via {DeepONet} based on the universal approximation theorem of operators.
\newblock \emph{Nature Machine Intelligence}, 3:218--229, 2021.

\bibitem[OpenFOAM(2024)]{OpenFOAM}
{OpenCFD Ltd.}
\newblock {OpenFOAM}: The open source {CFD} toolbox.
\newblock \url{https://www.openfoam.com}, 2024.

\bibitem[Oseledets(2011)]{Oseledets2011}
I.V.~Oseledets.
\newblock Tensor-train decomposition.
\newblock \emph{SIAM Journal on Scientific Computing}, 33(5):2295--2317, 2011.

\bibitem[Peherstorfer and Willcox(2016)]{Peherstorfer2016}
B.~Peherstorfer and K.~Willcox.
\newblock Data-driven operator inference for nonintrusive projection-based model reduction.
\newblock \emph{Computer Methods in Applied Mechanics and Engineering}, 306:196--215, 2016.

\bibitem[Port{\'e}-Agel et~al.(2020)]{PorteAgel2020}
F.~Port{\'e}-Agel, M.~Bastankhah, and S.~Shamsoddin.
\newblock Wind-turbine and wind-farm flows: A review.
\newblock \emph{Boundary-Layer Meteorology}, 174:1--59, 2020.

\bibitem[Tabib et~al.(2024)]{Tabib2024}
M.~Tabib and others.
\newblock Parametric reduced-order modeling for wind turbine wake flows using graph convolutional autoencoders.
\newblock \emph{In preparation}, 2024.

\end{thebibliography}

\end{document}
